\section{I. 流体力学}

\subsection{I1. 守恒律}

\textbf{质量守恒：}
\[
\frac{\partial \rho}{\partial t} + \nabla \cdot (\rho \mathbf{u}) = 0
\]

\textbf{动量守恒：}
\[
\frac{\partial(\rho \mathbf{u})}{\partial t} + \nabla \cdot (\rho \mathbf{u} \otimes \mathbf{u}) = \nabla \cdot \boldsymbol{\sigma} + \rho \mathbf{g}
\]

\textbf{物质导数：}
\[
\frac{D}{Dt} = \frac{\partial}{\partial t} + \mathbf{u} \cdot \nabla
\]

\subsection{I2. 应力张量与粘性}

\[
\boldsymbol{\sigma} = -p\mathbf{I} + \boldsymbol{\tau}
\]

\textbf{牛顿流体本构关系：}
\[
\boldsymbol{\tau} = 2\mu\mathbf{D} + \lambda(\nabla\cdot\mathbf{u})\mathbf{I}
\]

其中 $\mathbf{D} = \frac{1}{2}(\nabla\mathbf{u} + (\nabla\mathbf{u})^T)$ 是应变率张量。

\textbf{非牛顿流体：}
\begin{center}
\begin{tabular}{lll}
\toprule
\textbf{类型} & \textbf{特征} & \textbf{例子} \\
\midrule
剪切稀化 & 粘度随剪切率增大而减小 & 油漆、血液 \\
剪切增稠 & 粘度随剪切率增大而增大 & 玉米淀粉水 \\
宾汉流体 & 有屈服应力 & 牙膏、泥浆 \\
粘弹性 & 同时有粘性和弹性 & 聚合物、面团 \\
\bottomrule
\end{tabular}
\end{center}

\subsection{I3. 雷诺数}

\[
Re = \frac{UL}{\nu} = \frac{\text{惯性力}}{\text{粘性力}}
\]

\begin{center}
\begin{tabular}{ll}
\toprule
\textbf{$Re$ 范围} & \textbf{流动特征} \\
\midrule
$Re \ll 1$ & 蠕动流（Stokes 流），粘性主导 \\
$Re \sim 10^2$--$10^3$ & 层流，可能出现分离 \\
$Re > 4000$（管流） & 湍流 \\
$Re \sim 10^6+$ & 完全湍流 \\
\bottomrule
\end{tabular}
\end{center}

\subsection{I4. 涡量动力学}

\textbf{涡量：} $\boldsymbol{\omega} = \nabla \times \mathbf{u}$

\textbf{涡量方程（不可压缩）：}
\[
\frac{D\boldsymbol{\omega}}{Dt} = (\boldsymbol{\omega}\cdot\nabla)\mathbf{u} + \nu\nabla^2\boldsymbol{\omega}
\]

$(\boldsymbol{\omega}\cdot\nabla)\mathbf{u}$：涡量拉伸项（三维特有，湍流能量级联的关键）。

\textbf{Kelvin 环量定理（无粘）：}
\[
\frac{D\Gamma}{Dt} = 0, \quad \Gamma = \oint_C \mathbf{u}\cdot d\mathbf{l}
\]

在 CG 中：涡量约束（Vorticity Confinement）补偿数值耗散对涡量的抹平。

\subsection{I5. 湍流}

\textbf{Kolmogorov 能量谱：} $E(k) \sim k^{-5/3}$（惯性子区）

\textbf{Kolmogorov 尺度：}
\[
\eta = \left(\frac{\nu^3}{\epsilon}\right)^{1/4}
\]

完全解析湍流需要网格数 $\sim Re^{9/4}$，对 $Re = 10^6$ 需要 $\sim 10^{13}$ 个网格点——完全不可行。

\textbf{RANS：} 将速度分解为均值和脉动 $\mathbf{u} = \bar{\mathbf{u}} + \mathbf{u}'$，引入雷诺应力 $-\rho\overline{u_i'u_j'}$，需要模型封闭（$k$-$\varepsilon$、$k$-$\omega$ 等）。

\textbf{LES：} 直接解析大尺度，对小尺度用亚格子模型。

\subsection{I6. 边界层}

\textbf{边界层厚度：} $\delta \sim \dfrac{L}{\sqrt{Re}}$

\textbf{无滑移条件：} $\mathbf{u}|_{\text{wall}} = 0$

\subsection{I7. 可压缩流与激波}

\textbf{马赫数：} $Ma = \dfrac{U}{c}$

$Ma > 1$ 时出现激波。Rankine--Hugoniot 条件描述激波两侧的守恒关系。

\subsection{I8. 自由表面与多相流}

\begin{itemize}
  \item \textbf{Level Set：} 用符号距离函数 $\phi$ 表示界面（$\phi = 0$ 为界面），输运方程 $\frac{\partial \phi}{\partial t} + \mathbf{u}\cdot\nabla\phi = 0$
  \item \textbf{VOF：} 用体积分数 $f \in [0,1]$ 表示，天然守恒
  \item \textbf{SPH：} 拉格朗日描述，天然处理自由表面
\end{itemize}

\textbf{表面张力：} $\Delta p = \sigma \kappa$

\subsection{I9. 数值模拟中的守恒性}

\subsubsection{有限体积法——天然守恒}

相邻单元共享面的通量大小相等方向相反 $\to$ 全局守恒自动满足。

\subsubsection{能量稳定格式}

连续情况：
\[
\frac{d}{dt}\int \frac{1}{2}|\mathbf{u}|^2\,dV = -\nu\int|\nabla\mathbf{u}|^2\,dV \leq 0
\]

离散要求：对流项的离散必须满足"反对称性"或"守恒性"。

\subsubsection{散度自由条件的保持}

\begin{itemize}
  \item 投影法（Chorin/Temam）
  \item MAC/Staggered 网格：离散散度精确为零
  \item 约束输运：对涡量或流函数求解
\end{itemize}

\subsubsection{时间步长约束}

\begin{center}
\begin{tabular}{ll}
\toprule
\textbf{约束来源} & \textbf{条件} \\
\midrule
对流（CFL） & $\Delta t \leq C\frac{\Delta x}{|\mathbf{u}|_{\max}}$ \\
粘性 & $\Delta t \leq C\frac{\Delta x^2}{\nu}$ \\
表面张力 & $\Delta t \leq C\sqrt{\frac{\rho\,\Delta x^3}{\sigma}}$ \\
弹性波 & $\Delta t \leq C\frac{\Delta x}{c_s}$ \\
\bottomrule
\end{tabular}
\end{center}